%%%%%
%%%%%  Naudokite LUALATEX, ne LATEX.
%%%%%
%%%%
\documentclass[]{VUMIFTemplateClass}

\usepackage{indentfirst}
\usepackage{amsmath, amsthm, amssymb, amsfonts}
\usepackage{mathtools}
\usepackage{physics}
\usepackage{graphicx}
\usepackage{verbatim}
\usepackage[hidelinks]{hyperref}
\usepackage{color,algorithm,algorithmic}
\usepackage[nottoc]{tocbibind}
\usepackage{tocloft}
\usepackage[version=4]{mhchem}
\usepackage{textgreek}
\usepackage{textcomp}
\usepackage{float}
\usepackage{booktabs}

\usepackage{titlesec}
\newcommand{\sectionbreak}{\clearpage}

\makeatletter
\renewcommand{\fnum@algorithm}{\thealgorithm}
\makeatother
\renewcommand\thealgorithm{\arabic{algorithm} algoritmas}

\usepackage{biblatex}
\bibliography{bibliografija}

\newcommand{\EE}{\mathbb{E}\,}
\newcommand{\ee}{{\mathrm e}}
\newcommand{\RR}{\mathbb{R}}

\studijuprograma{Informatikos}
\darbotipas{Magistro baigiamojo darbo I dalis}
\darbopavadinimas{Duomenų paruošimas ir K-NN prognozavimas}
\darbopavadinimasantras{Data preparation and K-NN forecasting}
\autorius{Mindaugas Žiemys}

\vadovas{Linas Litvinas, Asist., Dr.}

\begin{document}
    \selectlanguage{lithuanian}

    \onehalfspacing
    \input{titulinis}

    \singlespacing

    \tableofcontents
    \onehalfspacing


    \section{Duomenų paruošimas}

    \subsection{Duomenų valymas}

    Duomenų rinkinys apima 2012--2025 metų valandinius matavimus iš Londono Merilebono kelio stoties -- iš viso 122\,736 eilutės.
    Iš jų 30\,770 eilučių (25,1\,\%) turi bent vieną trūkstamą reikšmę. \ref{tab:missing} lentelėje pateikiamas trūkstamų reikšmių skaičius ir procentas kiekvienam požymiui.

    \begin{table}[H]
        \centering
        \caption{Trūkstamos reikšmės pagal požymį}
        \small
        \setlength{\tabcolsep}{6pt}
        \begin{tabular}{|l|r|r|}
            \hline
            \textbf{Požymis} & \textbf{Trūkstamų reikšmių} & \textbf{\%} \\
            \hline
            \ce{CO}          & 12\,704                     & 10,4        \\
            \ce{KD_{10}}     & 10\,379                     & 8,5         \\
            \ce{KD_{2,5}}    & 9\,649                      & 7,9         \\
            \ce{O3}          & 4\,359                      & 3,6         \\
            \ce{NO}          & 3\,653                      & 3,0         \\
            \ce{NO2}         & 3\,653                      & 3,0         \\
            Vėjo greitis     & 3\,168                      & 2,6         \\
            Temperatūra      & 3\,168                      & 2,6         \\
            Vėjo kryptis     & 3\,168                      & 2,6         \\
            \hline
        \end{tabular}
        \label{tab:missing}
    \end{table}

    \ce{CO} turi didžiausią trūkstamų reikšmių dalį iš visų požymių (10,4\,\%). Dėl šių priežasčių \ce{CO} pašalintas iš duomenų rinkinio.
    Pašalinus \ce{CO}, eilučių su bent viena trūkstama reikšme liko 21\,636 (17,6\,\%).

    Trūkstamų reikšmių analizė parodo, kad 66,9\% visų trūkstamų reikšmių sudaro tarpai, ilgesni nei 24 val. \ref{tab:gap_distribution} lentelėje pateikiamas trūkstamų reikšmių pasiskirstymas pagal tarpo ilgį.

    \begin{table}[H]
        \centering
        \caption{Trūkstamų reikšmių pasiskirstymas pagal tarpo ilgį}
        \small
        \setlength{\tabcolsep}{5pt}
        \begin{tabular}{|l|r|r|r|}
            \hline
            \textbf{Tarpo ilgis} & \textbf{NaN reikšmių} & \textbf{\% visų NaN} \\
            \hline
            1 val.               & 1\,457                & 3,3                  \\
            2 val.               & 838                   & 1,9                  \\
            3--5 val.            & 1\,085                & 2,4                  \\
            6--12 val.           & 1\,059                & 2,4                  \\
            12--24 val.          & 10\,230               & 23,1                 \\
            24+ val.             & 29\,696               & 66,9                 \\
            \hline
            \textbf{Iš viso}     & \textbf{44\,365}      & \textbf{100\,\%}     \\
            \hline
        \end{tabular}
        \label{tab:gap_distribution}
    \end{table}

    Interpoliavus tarpus iki 5 val. prisidėtų tik $+$1\,301 papildoma eilutė (1\% duomenų rinkinio).
    Visos trūkstamos eilutės yra tiesiog pašalinamos.

    \subsection{Duomenų savybių inžinerija}

    \paragraph{Vėjo kryptis.}
    Vėjo kryptis (laipsniais, 0--360\textdegree) konvertuota į dvi dedamąsias: sinusą ir kosinusą:

    \begin{equation}
        s = \sin\!\left(\theta \cdot \frac{\pi}{180}\right), \qquad
        c = \cos\!\left(\theta \cdot \frac{\pi}{180}\right)
        \label{eq:windsincos}
    \end{equation}

    Toks transformavimas išsprendžia cikliškumo problemą, kampai 0\textdegree{} ir 360\textdegree{} yra identiški, tačiau skaitiškai skiriasi: $|360 - 0| = 360$, nors faktinis skirtumas 0\textdegree.
    Sinuso ir kosinuso atvaizdavime šie kampai duoda identiškas reikšmių poras.
    Originalus vėjo krypties stulpelis pašalintas.
    Sinusas ir kosinusas jau patenka į $[-1, 1]$ intervalą, todėl papildomo normalizavimo nereikia.

    \paragraph{Dienos kategorija.}
    Kiekviena laiko eilutė susieta su dienos kategorija, kurią galima aprašyti:

    \begin{equation}
        \text{kategorija}(d) =
        \begin{cases}
            1, & \text{jei } d \text{ yra poilsio diena} \\
            2, & \text{jei } d+1 \text{ yra poilsio diena} \\
            0, & \text{kitais atvejais (darbo diena)}
        \end{cases}
        \label{eq:daycat}
    \end{equation}

    čia poilsio diena -- savaitgalis (šeštadienis arba sekmadienis) arba JK valstybinė šventė. JK švenčių sąrašas generuojamas automatiškai \textit{Python} bibliotekos \texttt{holidays} pagalba, apimant 2012--2025 metus.

    \subsection{Duomenų statistika}

    \ref{tab:stats} lentelėje pateikiama duomenų statistika visiems normalizuojamiems požymiams, išmetus CO stulpelį bei eilutes su trūkstamomis reikšmėmis.

    \begin{table}[H]
        \centering
        \caption{Duomenų statistika po valymo}
        \small
        \setlength{\tabcolsep}{4pt}
        \begin{tabular}{|l|c|r|r|r|r|r|}
            \hline
            \textbf{Požymis} & \textbf{Vnt.}                  & \textbf{N} & \textbf{Vidurkis} & \textbf{Mediana} & \textbf{Min} & \textbf{Max} \\
            \hline
            \ce{O3}          & \textmu g/m\textsuperscript{3} & 101\,100   & 22,85             & 18,16            & $-0{,}93$    & 156,46       \\
            \ce{NO}          & \textmu g/m\textsuperscript{3} & 101\,100   & 85,93             & 48,43            & $-0{,}37$    & 872,83       \\
            \ce{NO2}         & \textmu g/m\textsuperscript{3} & 101\,100   & 66,40             & 58,52            & 0,00         & 321,91       \\
            \ce{KD_{10}}     & \textmu g/m\textsuperscript{3} & 101\,100   & 22,87             & 20,20            & $-2{,}90$    & 187,90       \\
            \ce{KD_{2,5}}    & \textmu g/m\textsuperscript{3} & 101\,100   & 14,51             & 11,94            & $-5{,}00$    & 127,60       \\
            Vėjo greitis     & m/s                            & 101\,100   & 3,46              & 3,20             & 0,00         & 13,70        \\
            Temperatūra      & \textdegree C                  & 101\,100   & 10,23             & 10,20            & $-10{,}50$   & 33,40        \\
            \hline
        \end{tabular}
        \label{tab:stats}
    \end{table}

    \ce{NO} pasižymi labai dideliu maksimumu (max 872,83 \textmu g/m\textsuperscript{3}), 18 kartų didesniu nei mediana. \ce{KD_{2,5}} maksimumas (127,60 \textmu g/m\textsuperscript{3}) daugiau nei 8 kartus didesnis už medianą. Tai rodo retų, bet intensyvių taršos epizodų egzistavimą.

    \subsection{Sezoniniai teršalų trendai}

    \ref{tab:seasonal} lentelėje pateikiami pagrindinių teršalų mėnesiniai vidurkiai.
    Sausio -- kovo mėn. \ce{KD_{2,5}} padidėja. NO padidėja lapkričio -- sausio mėn.

    \begin{table}[H]
        \centering
        \caption{Teršalų mėnesiniai vidurkiai}
        \small
        \setlength{\tabcolsep}{5pt}
        \begin{tabular}{|l|r|r|r|r|r|}
            \hline
            \textbf{Mėnuo} & \textbf{\ce{O3}} & \textbf{\ce{NO}} & \textbf{\ce{NO2}} & \textbf{\ce{KD_{10}}} & \textbf{\ce{KD_{2,5}}} \\
            \hline
            Sausis         & 15,82            & 108,35           & 70,77             & 25,35                 & 16,23                  \\
            Vasaris        & 19,97            & 96,44            & 69,10             & 25,75                 & 16,12                  \\
            Kovas          & 25,85            & 80,21            & 68,79             & 27,86                 & 18,62                  \\
            Balandis       & 34,58            & 67,69            & 65,29             & 24,27                 & 15,59                  \\
            Gegužė         & 34,50            & 64,23            & 62,18             & 21,43                 & 13,68                  \\
            Birželis       & 29,32            & 72,99            & 64,67             & 19,54                 & 12,27                  \\
            Liepa          & 24,30            & 74,79            & 64,86             & 18,40                 & 11,87                  \\
            Rugpjūtis      & 23,17            & 69,61            & 60,90             & 20,55                 & 12,07                  \\
            Rugsėjis       & 19,59            & 81,90            & 64,56             & 21,66                 & 13,84                  \\
            Spalis         & 15,14            & 95,07            & 66,21             & 22,73                 & 13,64                  \\
            Lapkritis      & 13,90            & 106,25           & 66,16             & 23,94                 & 15,08                  \\
            Gruodis        & 16,71            & 106,22           & 67,81             & 22,46                 & 14,55                  \\
            \hline
        \end{tabular}
        \label{tab:seasonal}
    \end{table}

    Galima būtų tobulinti K-NN metodą: kaimynų paiešką galima apriboti tuo pačiu sezonu arba mėnesiu, arba sezoninius faktorius įtraukti kaip papildomą svorį Euklido atstumo formulėje.

    \subsection{Min-max normalizavimas}

    Skaitiniai matavimų atributai normalizuoti į $[-1, 1]$ intervalą pagal formulę:

    \begin{equation}
        x' = \frac{2(x - x_{\min})}{x_{\max} - x_{\min}} - 1
        \label{eq:minmax}
    \end{equation}

    čia $x$ -- originali reikšmė, $x_{\min}$ ir $x_{\max}$ -- atitinkamo atributo minimali ir maksimali reikšmės visoje duomenų aibėje, $x'$ -- normalizuota reikšmė. \ref{tab:minmax} lentelėje pateikiamos normalizuojamų atributų minimalios ir maksimalios reikšmės.

    \begin{table}[H]
        \centering
        \caption{Normalizuojamų atributų minimalios ir maksimalios reikšmės}
        \small
        \setlength{\tabcolsep}{5pt}
        \begin{tabular}{|l|c|r|r|}
            \hline
            \textbf{Atributas} & \textbf{Matavimo vienetas}     & \textbf{$x_{\min}$} & \textbf{$x_{\max}$} \\
            \hline
            \ce{O3}            & \textmu g/m\textsuperscript{3} & $-0{,}93$           & $156{,}46$          \\
            \ce{NO}            & \textmu g/m\textsuperscript{3} & $-0{,}37$           & $872{,}83$          \\
            \ce{NO2}           & \textmu g/m\textsuperscript{3} & $0{,}00$            & $321{,}91$          \\
            \ce{KD_{10}}       & \textmu g/m\textsuperscript{3} & $-2{,}90$           & $187{,}90$          \\
            \ce{KD_{2,5}}      & \textmu g/m\textsuperscript{3} & $-5{,}00$           & $127{,}60$          \\
            Vėjo greitis       & m/s                            & $0{,}00$            & $13{,}70$           \\
            Temperatūra        & \textdegree C                  & $-10{,}50$          & $33{,}40$           \\
            \hline
        \end{tabular}
        \label{tab:minmax}
    \end{table}

    Vėjo krypties sinuso ir kosinuso komponentai nebuvo normalizuoti papildomai, nes pagal apibrėžimą jau patenka į $[-1,\,1]$ intervalą.
    Dienos kategorija ir valanda yra kategoriniai požymiai, jų normalizavimas būtų beprasmis.

    \printbibliography[title = {Literatūra ir šaltiniai}]

\end{document}
